% book12.tex --- Book XII of Euclid's Elements: Method of Exhaustion.
%
% All 18 propositions encoded.  Book XII develops the Eudoxean method of
% exhaustion (X.1) to compute curvilinear and 3D volumes: circles
% (XII.2), pyramids (XII.5-XII.9), cones (XII.10-XII.15), and spheres
% (XII.16-XII.18).
%
% Wording follows Heath (1908).

\section{Book XII --- Method of Exhaustion}
\label{sec:book-XII}

\begin{claim}[Proposition XII.1: Similar inscribed polygons in circles are as squares on diameters]
\label{prop:XII.1}
Similar polygons inscribed in circles are to one another as the
squares on the diameters.
\end{claim}
\begin{evidence}[Proof of XII.1]
\label{ev:XII.1}
Decompose each polygon into similar triangles by joining vertices to
the centres; each pair of corresponding triangles is similar (VI.20)
with side ratio equal to the diameter ratio; sum and combine via
V.12.
\dependson{XII.1}{VI.20}
\dependson{XII.1}{V.12}
\end{evidence}

\begin{claim}[Proposition XII.2: Circles are as squares on diameters]
\label{prop:XII.2}
Circles are to one another as the squares on the diameters.
\end{claim}
\begin{evidence}[Proof of XII.2]
\label{ev:XII.2}
Apply the method of exhaustion: inscribe similar polygons in the two
circles; by XII.1 they are in the ratio of squares on the diameters.
Any deviation from that ratio at the level of the circles leads, via
X.1, to a contradiction by choosing inscribed polygons close enough
to fill the circle.
\dependson{XII.2}{X.1}
\dependson{XII.2}{XII.1}
\end{evidence}

\begin{claim}[Proposition XII.3: A pyramid on a triangular base is split into two equal smaller pyramids and two equal prisms]
\label{prop:XII.3}
Any pyramid which has a triangular base is divided into two pyramids
equal and similar to one another, similar to the whole, and having
triangular bases, and into two equal prisms; and the two prisms are
greater than the half of the whole pyramid.
\end{claim}
\begin{evidence}[Proof of XII.3]
\label{ev:XII.3}
Bisect each edge (I.10); the midpoint cuts split the pyramid into
two corner pyramids and two prisms.  The two corner pyramids are
similar to the original (their edges halved), and by I.34 the two
prisms are congruent.
\dependson{XII.3}{I.10}
\dependson{XII.3}{I.34}
\dependson{XII.3}{XI.39}
\end{evidence}

\begin{claim}[Proposition XII.4: Pyramids of equal height are as their bases (special case)]
\label{prop:XII.4}
If there be two pyramids of the same height which have triangular
bases, and each of them be divided into two pyramids equal to one
another and similar to the whole, and into two equal prisms, then,
as the base of the one pyramid is to the base of the other pyramid,
so will all the prisms in the one pyramid be to all the prisms in
the other pyramid.
\end{claim}
\begin{evidence}[Proof of XII.4]
\label{ev:XII.4}
Each iteration of XII.3 doubles the number of prisms; by
proportionality of bases (VI.1) the prism-sums maintain the same
ratio as the original bases.
\dependson{XII.4}{VI.1}
\dependson{XII.4}{XII.3}
\end{evidence}

\begin{claim}[Proposition XII.5: Pyramids of equal height are as their bases]
\label{prop:XII.5}
Pyramids which are of the same height and have triangular bases are
to one another as their bases.
\end{claim}
\begin{evidence}[Proof of XII.5]
\label{ev:XII.5}
Apply XII.4 in the limit of XII.3 iterations; the prism-sums
exhaust the pyramids (X.1), so the base-ratio is the pyramid-ratio.
\dependson{XII.5}{X.1}
\dependson{XII.5}{XII.3}
\dependson{XII.5}{XII.4}
\end{evidence}

\begin{claim}[Proposition XII.6: Pyramids of equal height on polygonal bases are as their bases]
\label{prop:XII.6}
Pyramids which are of the same height and have polygonal bases are
to one another as the bases.
\end{claim}
\begin{evidence}[Proof of XII.6]
\label{ev:XII.6}
Triangulate each polygonal base; the polygon-pyramid is the sum of
the triangular sub-pyramids; apply XII.5 and V.12.
\dependson{XII.6}{V.12}
\dependson{XII.6}{XII.5}
\end{evidence}

\begin{claim}[Proposition XII.7: Triangular prism is three equal pyramids]
\label{prop:XII.7}
Any prism which has a triangular base is divided into three pyramids
equal to one another which have triangular bases.
\end{claim}
\begin{evidence}[Proof of XII.7]
\label{ev:XII.7}
Cut the prism by two planes through opposite edge-pairs; the three
resulting pyramids share a common apex and have congruent base
triangles, so by XII.5 they have equal volume.
\dependson{XII.7}{XI.39}
\dependson{XII.7}{XII.5}
\end{evidence}

\begin{claim}[Proposition XII.8: Similar pyramids are as the cubes on corresponding edges]
\label{prop:XII.8}
Similar pyramids which have triangular bases are in the triplicate
ratio of their corresponding sides.
\end{claim}
\begin{evidence}[Proof of XII.8]
\label{ev:XII.8}
By XII.7 a prism is three equal pyramids; by XI.33 similar
parallelepipeds (and hence prisms) are in the triplicate ratio of
edges; transfer to pyramids by XII.5.
\dependson{XII.8}{XI.33}
\dependson{XII.8}{XII.5}
\dependson{XII.8}{XII.7}
\end{evidence}

\begin{claim}[Proposition XII.9: Equal pyramids have reciprocally proportional bases and heights]
\label{prop:XII.9}
In equal pyramids which have triangular bases the bases are
reciprocally proportional to the heights; and those pyramids which
have triangular bases in which the bases are reciprocally
proportional to the heights are equal.
\end{claim}
\begin{evidence}[Proof of XII.9]
\label{ev:XII.9}
3D analogue of VI.15 for pyramids; via XII.5 / XII.6 the area-times-
height proportion factors into the reciprocal proportion of bases
and heights.
\dependson{XII.9}{VI.14}
\dependson{XII.9}{XII.5}
\dependson{XII.9}{XII.6}
\end{evidence}

\begin{claim}[Proposition XII.10: A cone is one-third of the cylinder on the same base and height]
\label{prop:XII.10}
Any cone is a third part of the cylinder which has the same base
with it and equal height.
\end{claim}
\begin{evidence}[Proof of XII.10]
\label{ev:XII.10}
Inscribe a pyramid on a polygonal base in both cone and cylinder;
XII.7 makes the pyramid one-third the prism; apply X.1 to refine the
inscribed polygon to fill the circle (XII.2); the limit gives the
cone-to-cylinder ratio.
\dependson{XII.10}{X.1}
\dependson{XII.10}{XII.2}
\dependson{XII.10}{XII.7}
\end{evidence}

\begin{claim}[Proposition XII.11: Cones and cylinders of equal height are as their bases]
\label{prop:XII.11}
Cones and cylinders which are of the same height are to one another
as their bases.
\end{claim}
\begin{evidence}[Proof of XII.11]
\label{ev:XII.11}
By XII.2 the bases (circles) are in the squared-diameter ratio; by
the formula in XII.10, the volumes follow the same ratio.
\dependson{XII.11}{XII.2}
\dependson{XII.11}{XII.10}
\end{evidence}

\begin{claim}[Proposition XII.12: Similar cones and cylinders are as cubes on diameters]
\label{prop:XII.12}
Similar cones and cylinders are to one another in the triplicate
ratio of the diameters in their bases.
\end{claim}
\begin{evidence}[Proof of XII.12]
\label{ev:XII.12}
Analogue of XII.8 for cones / cylinders: by similarity the height
scales proportionally with the diameter; cube of the linear ratio
gives the volume ratio.
\dependson{XII.12}{XII.8}
\dependson{XII.12}{XII.10}
\dependson{XII.12}{XII.11}
\end{evidence}

\begin{claim}[Proposition XII.13: Parallel sections of a cylinder are in ratio of distances]
\label{prop:XII.13}
If a cylinder be cut by a plane which is parallel to its opposite
planes, then, as the cylinder is to the cylinder, so will the axis
be to the axis.
\end{claim}
\begin{evidence}[Proof of XII.13]
\label{ev:XII.13}
Parallel cross-sections give equal circles (XI.16 implies parallel
diameters); the volume scales linearly with axial length by XII.11.
\dependson{XII.13}{XI.16}
\dependson{XII.13}{XII.11}
\end{evidence}

\begin{claim}[Proposition XII.14: Cylinders on equal bases are as their heights]
\label{prop:XII.14}
Cones and cylinders which are on equal bases are to one another as
their heights.
\end{claim}
\begin{evidence}[Proof of XII.14]
\label{ev:XII.14}
By XII.13 the volume is proportional to the axis when the base is
fixed.
\dependson{XII.14}{XII.11}
\dependson{XII.14}{XII.13}
\end{evidence}

\begin{claim}[Proposition XII.15: Equal cones / cylinders have reciprocal proportions]
\label{prop:XII.15}
In equal cones and cylinders the bases are reciprocally proportional
to the heights; and those cones and cylinders in which the bases
are reciprocally proportional to the heights are equal.
\end{claim}
\begin{evidence}[Proof of XII.15]
\label{ev:XII.15}
Analogue of XII.9 / VI.15 for cones and cylinders.
\dependson{XII.15}{XII.9}
\dependson{XII.15}{XII.11}
\dependson{XII.15}{XII.14}
\end{evidence}

\begin{claim}[Proposition XII.16: Inscribe in the larger of two concentric circles a polygon not touching the smaller]
\label{prop:XII.16}
Given two circles about the same centre, to inscribe in the greater
circle an equilateral polygon with an even number of sides which
does not touch the lesser circle.
\end{claim}
\begin{evidence}[Proof of XII.16]
\label{ev:XII.16}
Bisect arcs repeatedly (III.30) until the chord-to-arc gap is
smaller than the difference of radii; this guarantees that the
inscribed polygon avoids touching the smaller circle.
\dependson{XII.16}{III.30}
\dependson{XII.16}{X.1}
\end{evidence}

\begin{claim}[Proposition XII.17: Inscribe in the larger of two concentric spheres a polyhedron not touching the smaller]
\label{prop:XII.17}
Given two spheres about the same centre, to inscribe in the greater
sphere a polyhedral solid which does not touch the lesser sphere at
its surface.
\end{claim}
\begin{evidence}[Proof of XII.17]
\label{ev:XII.17}
3D analogue of XII.16: apply XII.16 in great-circle cross-sections,
then triangulate the sphere using XI.27 to assemble a polyhedron
strictly inside the larger sphere and outside the smaller.
\dependson{XII.17}{XI.27}
\dependson{XII.17}{XII.16}
\end{evidence}

\begin{claim}[Proposition XII.18: Spheres are in triplicate ratio of diameters]
\label{prop:XII.18}
Spheres are to one another in the triplicate ratio of their
respective diameters.
\end{claim}
\begin{evidence}[Proof of XII.18]
\label{ev:XII.18}
Apply the method of exhaustion: inscribe similar polyhedra (XII.17);
by XII.12 (similar cones) and the polyhedron's similar-pyramid
decomposition, the inscribed solids are in the triplicate ratio of
diameters.  By X.1 the inscribed solids approach the spheres in
volume; the limit gives the result.
\dependson{XII.18}{X.1}
\dependson{XII.18}{XII.8}
\dependson{XII.18}{XII.12}
\dependson{XII.18}{XII.17}
\end{evidence}
